\chapter{\abstractname}

This thesis presents the design, modeling, and experimental validation of an X-wing hybrid aerial platform that combines quadrotor vertical take-off and landing (VTOL) capability with fixed-wing aerodynamic efficiency. The platform features four wings arranged in an X-configuration, providing passive lift that reduces required rotor thrust during forward flight.

A comprehensive dynamics model incorporating aerodynamic forces and moments is developed and integrated with an Incremental Nonlinear Dynamic Inversion (INDI) control architecture. The controller requires no explicit aerodynamic model and treats wing-generated forces as bounded disturbances, enabling robust trajectory tracking across the entire flight envelope from hover to high-speed cruise.

Experimental validation was conducted in a motion capture arena with a \SI{2.5}{\kilogram} prototype. Two configurations were tested: NACA~0015 airfoils at \SI{45}{\degree} dihedral and AG25 airfoils at \SI{20}{\degree} dihedral. Thrust characterization yielded the mapping $T = 5.15 \times 10^{-6} \cdot \text{RPM}^2 - 0.090 \times 10^{-2} \cdot \text{RPM}$ (N) with $R^2 = 0.997$. Flight tests demonstrated stable, accurate tracking up to \SI{8}{\meter\per\second} with RMS position errors below \SI{0.24}{\meter}, sustaining centripetal accelerations up to 2.8g, demonstrating stable control under strong unmodeled aerodynamic loads and indicative rotor unloading.

Electrical power was directly measured in outdoor flight using synchronized current and voltage telemetry. Results show a measurable reduction in mechanical power demand with increasing airspeed—up to 33\% at 10\si{\meter\per\second}—confirming that aerodynamic lift from the fixed surfaces translates to tangible energy savings while maintaining full control authority. 

\chapter*{Zusammenfassung}
\addcontentsline{toc}{chapter}{Zusammenfassung}

Diese Arbeit präsentiert den Entwurf, die Modellierung und die experimentelle Validierung einer X-Flügel-Hybrid-Luftplattform, die die vertikale Start- und Landefähigkeit (VTOL) eines Quadrocopters mit der aerodynamischen Effizienz eines Starrflüglers kombiniert. Die Plattform verfügt über vier Flügel, die in einer X-Konfiguration angeordnet sind und einen passiven Auftrieb erzeugen, der den erforderlichen Rotorschub während des Vorwärtsflugs reduziert.

Es wurde ein umfassendes Dynamikmodell entwickelt, das aerodynamische Kräfte und Momente berücksichtigt und in eine INDI-Steuerungsarchitektur (Incremental Nonlinear Dynamic Inversion) integriert wurde. Der Regler benötigt kein explizites aerodynamisches Modell und behandelt die von den Flügeln erzeugten Kräfte als begrenzte Störungen, was eine robuste Flugbahnverfolgung über den gesamten Flugbereich vom Schwebeflug bis zum Hochgeschwindigkeitsflug ermöglicht.

Die experimentelle Validierung wurde in einer Motion-Capture-Arena mit einem \SI{2,5}{\kilogram} schweren Prototyp durchgeführt. Es wurden zwei Konfigurationen getestet: NACA~0015-Tragflächen mit \SI{45}{\degree} V-Form und AG25-Tragflächen mit \SI{20}{\degree} V-Form. Die Schubcharakterisierung ergab die Zuordnung $T = 5{,}15 \times 10^{-6} \cdot \text{RPM}^2 - 0{,}090 \times 10^{-2} \cdot \text{RPM}$ (N) mit $R^2 = 0{,}997$. Flugtests zeigten eine stabile, genaue Verfolgung bis zu \SI{8}{\meter\per\second} mit RMS-Positionsfehlern unter \SI{0,24}{\meter}, bei anhaltenden Zentripetalbeschleunigungen bis zu 2,8g, was eine stabile Steuerung unter starken, nicht modellierten aerodynamischen Belastungen und einer indikativen Rotorentlastung demonstriert.

Die elektrische Leistung wurde im Außenflug direkt mit synchronisierter Strom- und Spannungstelemetrie gemessen. Die Ergebnisse zeigen eine messbare Verringerung des mechanischen Leistungsbedarfs mit zunehmender Fluggeschwindigkeit -- bis zu 33\% bei \SI{10}{\meter\per\second} -- was bestätigt, dass der aerodynamische Auftrieb der festen Flächen zu spürbaren Energieeinsparungen führt, während die volle Regelautorität erhalten bleibt.
